% titel: Blatt 6, A2: Rechnen mit komplexen Zahlen
% Teil II: der durchgerechnete Weg.

\subsection*{Gegeben und gesucht}
Lösungsmenge von $|z-1+i|>1$ (a), alle $z$ mit $|z+i|=1$ und $|z-i|=\sqrt3$
(b), Nachweis von $|zw|=|z||w|$ für alle $z,w\in\mathbb C$ (c).

\subsection*{a) $|z-1+i|>1$}
\begin{equation*}
  |z-1+i|=|z-(1-i)|.
\end{equation*}
Das ist der Abstand von $z$ zum Punkt $1-i$. Die Ungleichung $>1$ beschreibt
das \textbf{komplementäre Gebiet außerhalb der offenen Kreisscheibe} um
$1-i$ mit Radius $1$ (ohne Rand).

\subsection*{b) Schnitt zweier Kreise}
Sei $z=x+iy$. Die Bedingungen quadrieren:
\begin{align*}
  |z+i|^2&=x^2+(y+1)^2=1,\\
  |z-i|^2&=x^2+(y-1)^2=3.
\end{align*}
Subtraktion der ersten von der zweiten Zeile:
\begin{equation*}
  x^2+(y-1)^2-\bigl[x^2+(y+1)^2\bigr]=3-1
  \;\Longrightarrow\;
  (y-1)^2-(y+1)^2=2
  \;\Longrightarrow\;
  -4y=2,\quad y=-\tfrac12.
\end{equation*}
Einsetzen in $x^2+(y+1)^2=1$ mit $y+1=\tfrac12$:
\begin{equation*}
  x^2+\tfrac14=1 \;\Longrightarrow\; x^2=\tfrac34
  \;\Longrightarrow\; x=\pm\tfrac{\sqrt3}{2}.
\end{equation*}
Es gibt \textbf{zwei} Lösungen: $z=\pm\frac{\sqrt3}{2}-\frac12i$.

\subsection*{c) $|zw|=|z||w|$}
Mit $|u|^2=u\bar u$:
\begin{equation*}
  |zw|^2=(zw)\overline{(zw)}=zw\,\bar z\,\bar w=(z\bar z)(w\bar w)
  =|z|^2\,|w|^2.
\end{equation*}
Da der Betrag nichtnegativ ist, folgt $|zw|=|z||w|$.

\begin{ergebnis}
\begin{equation*}
  \text{a)}\ |z-(1-i)|>1\ \text{(äußerer Bereich um }1-i); \qquad
  \text{b)}\ z=\pm\tfrac{\sqrt3}{2}-\tfrac12i; \qquad
  \text{c)}\ |zw|=|z||w|.
\end{equation*}
\end{ergebnis}

\begin{kontrolle}
\textbf{Probe b):} Für $z=\frac{\sqrt3}{2}-\frac12i$ ist
$|z+i|=|\frac{\sqrt3}{2}+\frac12i|=\sqrt{\frac34+\frac14}=1$; deutlich
$|z-i|=|\frac{\sqrt3}{2}-\frac32i|=\sqrt{\frac34+\frac94}=\sqrt3$. Beide
Bedingungen erfüllt, ebenso für das Minus-Vorzeichen.

\textbf{Beispiel c):} $z=i$, $w=1+i$: $|zw|=|i-1|=\sqrt2$,
$|i||1+i|=1\cdot\sqrt2$ -- passt.
\end{kontrolle}

\begin{merksatz}
$|z-z_0|$ zeichnet sich als Kreis; Ungleichungen mit $<$/$>$ trennen Innen
und Außen. Schnitt zweier Betragsbedingungen: quadrieren, subtrahieren
(dadurch linear), auflösen -- jeder Schnittkreispartner gibt eine Lösung.
$|zw|=|z||w|$ folgt aus $|u|^2=u\bar u$.
\end{merksatz}
